\documentclass[conference]{IEEEtran}
\IEEEoverridecommandlockouts
% The preceding line is only needed to identify funding in the first footnote. If that is unneeded, please comment it out.
\usepackage{cite}
\usepackage{stfloats}
\usepackage{amsmath,amssymb,amsfonts}
\usepackage{algorithmic}
\usepackage{graphicx}
\graphicspath{{figures/}{./}}
\usepackage{textcomp}
\usepackage{url}
\usepackage{orcidlink}
\hypersetup{hidelinks}
\usepackage{xcolor}
\usepackage{pifont}
\newcommand{\ket}[1]{\left|#1\right\rangle}
\newcommand{\bra}[1]{\left\langle#1\right|}
\def\BibTeX{{\rm B\kern-.05em{\sc i\kern-.025em b}\kern-.08em
    T\kern-.1667em\lower.7ex\hbox{E}\kern-.125emX}}
\begin{document}

\title{Noise-Aware Simulation and Statistical Eavesdropping Detection in BB84, Six-State, and E91 Quantum Key Distribution Protocols\\
}

\author{
\IEEEauthorblockN{1\textsuperscript{st} Vidhi \orcidlink{0009-0002-8289-3182}}
\IEEEauthorblockA{\textit{Faculty of Computer Science} \\
\textit{Algoma University}\\
Brampton, Canada \\
vidhi01@algomau.ca}
\and
\IEEEauthorblockN{2\textsuperscript{nd} Muhammad Minhajuddin \orcidlink{0009-0003-0305-9497}}
\IEEEauthorblockA{\textit{Faculty of Computer Science} \\
\textit{Algoma University}\\
Brampton, Canada \\
mminhajuddin@algomau.ca}
\\[0.5em] 
\IEEEauthorblockN{4\textsuperscript{th} Ajmery Sultana \orcidlink{0000-0002-8591-026X} }
\IEEEauthorblockA{\textit{Faculty of Computer Science} \\
\textit{Algoma University}\\
Brampton, Canada \\
ajmery.sultana@algomau.ca}
\and
\IEEEauthorblockN{3\textsuperscript{rd} S. M. Riazul Islam \orcidlink{0000-0003-2968-9561}}
\IEEEauthorblockA{\textit{Dept. of Computing Science} \\.
\textit{University of Aberdeen} \\
Aberdeen, United Kingdom\\
riazul.islam@abdn.ac.uk}
}

\maketitle

\begin{abstract} The expected arrival of large scale quantum computers
threatens the security of classical public key cryptosystems such as RSA and
elliptic curve cryptography, whose security rests on computational hardness
assumptions that Shor's algorithm can defeat. Quantum Key Distribution (QKD)
offers an alternative whose security is grounded in the laws of quantum
mechanics rather than computational difficulty, allowing two parties to detect
eavesdropping on the channel before a secret key is used. Prior simulation
studies tend to evaluate single protocols, report a narrow set of metrics, and
model an eavesdropper that attacks either every transmitted qubit or none. As a
result, the way different QKD protocols expose a partial eavesdropper, and the
cost of that detection in terms of usable key, remain poorly characterized. In
this work, we present a unified Qiskit based study of 3 protocols, the 2
basis BB84 protocol, the 3 basis six state protocol, and the entanglement
based E91 protocol, under matched ideal and noisy channel conditions. Beyond the
standard metrics of Quantum Bit Error Rate (QBER), CHSH inequality violation,
and secure key rate, we characterize eavesdropping detection statistically
through the false positive rate and false negative rate, and we extend this
analysis from full interception to partial intercept and resend, where the
attacker targets only a fraction of the transmission. This yields, for each
protocol, a tradeoff between detection performance and secure key rate as a
function of attack strength and noise. Statistical detection of this kind was
established for BB84 and recently extended to the six state protocol for prepare
and measure QKD. We adopt those thresholds and provide the first false positive
and false negative analysis of entanglement based CHSH detection, together with
the first comparison of statistical detection across prepare and measure and
entanglement based protocols within a single framework. The study establishes a
reproducible environment that can be extended toward larger key lengths and
real device noise. Across matched ideal channels, a partial attacker who
intercepts up to 30\% of the transmission evades detection in all
3 protocols while a positive secret key is still produced, and per session
detection becomes reliable only near an interception fraction of 1/2, by
which point the secure key rate has already collapsed. The six state protocol
closes this stealth window earliest and needs the fewest compared signals to
suppress false alarms, BB84 is intermediate, and the entanglement based detector
of E91 is the most evadable and the most fragile to channel noise, with false
alarms rising sharply once the honest error rate exceeds about 8\%.
Under the measured \texttt{ibm\_marrakesh} device noise, however, all 3
protocols operate honestly without false alarms.
\end{abstract}

\begin{IEEEkeywords}
 Quantum Key Distribution, BB84, six-state protocol, E91, Qiskit, QBER, CHSH, secure key rate, eavesdropper detection, false positive rate, false negative rate
\end{IEEEkeywords}

\section{Introduction}

Cryptographic systems have long protected sensitive communication by encoding
plaintext into ciphertext that can only be recovered with the correct key.
Classical cryptography achieves confidentiality, authenticity, and integrity
through symmetric and asymmetric schemes such as the Advanced Encryption
Standard (AES), the Data Encryption Standard (DES), the Rivest Shamir Adleman
(RSA) algorithm, and Elliptic Curve Cryptography (ECC). The security of these
schemes depends on the computational difficulty of mathematical problems such
as integer factorization and the discrete logarithm. For decades, this
assumption has been sufficient, because no classical algorithm can solve these
problems efficiently at the key sizes used in practice.

The development of quantum computing changes this picture. Algorithms such as
Shor's algorithm and Grover's algorithm can solve, or substantially accelerate
the solution of, the problems on which classical public key cryptography
depends. A sufficiently large quantum computer would therefore render much of
the currently deployed asymmetric cryptography insecure. This threat is not
limited to future communication. In a harvest now, decrypt later attack, an
adversary records encrypted traffic today and decrypts it once a capable
quantum computer becomes available, which exposes data that is expected to
remain confidential for many years. This prospect has motivated an urgent
search for communication methods that remain secure against quantum
adversaries.

Quantum Key Distribution (QKD) addresses this problem from a fundamentally
different direction. Rather than relying on computational hardness, QKD derives
its security from the laws of quantum mechanics, in particular the no cloning
theorem and the measurement disturbance principle. Because an unknown quantum
state cannot be copied, and because any measurement of a quantum state in the
wrong basis disturbs it, an eavesdropper who attempts to intercept the quantum
channel inevitably introduces detectable errors. The legitimate parties,
traditionally called Alice and Bob, can therefore estimate the error rate
across a portion of their exchanged bits and decide whether the channel is
secure before deriving a final key. QKD does not transmit the message itself
over the quantum channel. Instead, it establishes a shared secret key that is
later used with a conventional encryption scheme, while an authenticated
classical channel supports basis reconciliation and error checking.

To make the mechanism of QKD concrete, it is useful to recall how quantum
information is represented and manipulated. The basic unit is the quantum bit,
or qubit, whose state is a unit vector in a 2 dimensional complex space
spanned by the computational basis states $\ket{0}$ and $\ket{1}$. Unlike a
classical bit, a qubit may occupy a superposition $\ket{\psi} = \alpha \ket{0} +
\beta \ket{1}$, with $|\alpha|^2 + |\beta|^2 = 1$. Measuring it in the
computational basis returns $0$ with probability $|\alpha|^2$ or $1$ with
probability $|\beta|^2$, and collapses the state onto whichever outcome occurs.
Single qubit gates rotate this state, and 3 of them suffice for the
protocols studied here. The Pauli $X$ gate exchanges $\ket{0}$ and $\ket{1}$,
acting as a bit flip. The Hadamard gate $H$ maps the computational states to the
equal superpositions $\ket{+}$ and $\ket{-}$, moving between the two conjugate
bases that BB84 uses. The phase gate $S$ adds a relative phase, which opens up a
third measurement basis. These bases are what make the protocols work: a qubit
gives a deterministic result when measured in the basis it was prepared in, but
a uniformly random result in a conjugate basis, so an eavesdropper who guesses
the wrong basis turns her mistake into observable errors. Two qubit gates such
as the controlled NOT create entanglement, a joint state that cannot be factored
into separate single qubit states and whose measurement outcomes stay correlated
no matter how far apart the 2 parties are. Encoding bits into single qubit
states, measuring them in randomly chosen bases, and testing the correlations of
entangled pairs are the operations the three protocols are built from, and a
circuit framework such as Qiskit expresses each one as a gate applied to a
simulated qubit.

The first and most widely studied QKD protocol, BB84, was introduced by Bennett
and Brassard in 1984 \cite{bennett1984quantum}. It encodes classical bits into
the polarization states of single photons using 2 conjugate bases, the
rectilinear basis at 0 and 90 degrees and the diagonal basis at 45
and 135 degrees. Alice prepares each qubit in a
randomly chosen basis, Bob measures each qubit in a randomly chosen basis, and
the two parties keep only the bits for which their bases agree, a step known as
sifting or basis reconciliation. The six state protocol generalizes BB84 by
adding a third mutually unbiased basis, so that Alice and Bob choose among
3 bases rather than 2. The additional basis lowers the chance that an
eavesdropper guesses the correct basis. This raises the disturbance an intercept
and resend attack imposes, from 25\% to 33\%, and
so improves the prospects for detection. The cost is a lower fraction of
matching bases, and therefore a shorter sifted key. A complementary
approach, the E91 protocol proposed by Ekert in 1991 \cite{ekert1991quantum},
replaces prepare and measure with shared entanglement. A source distributes
pairs of entangled qubits to Alice and Bob, who measure them along randomly
chosen axes. The correlations between their outcomes are tested against the
Clauser Horne Shimony Holt (CHSH) form of Bell's inequality, and a violation of
this inequality certifies that the entanglement was genuine and that no
eavesdropper has intercepted the qubits. These 3 protocols span the main
design choices in discrete variable QKD, namely the number of preparation bases
and the use of entanglement, which makes them a natural set for a controlled
comparison.

Although these protocols are well understood theoretically, evaluating their
behaviour under realistic conditions remains an active research problem.
Physical QKD hardware is expensive and inaccessible to most researchers, which
has made circuit level simulation an essential tool for studying protocol
performance. Open source quantum computing frameworks such as Qiskit allow
researchers to construct QKD circuits, introduce controlled noise through
configurable channel models, and measure key performance indicators on
classical hardware or cloud based simulators. However, existing simulation work
shares three limitations. First, protocols are usually studied in isolation or,
when compared, are reported only through coarse metrics such as the raw QBER.
Second, the statistical analysis of eavesdropping detection through the false
positive and false negative rates was first developed for BB84 and has recently
been extended to the six state protocol, but it has never been applied to
entanglement based detection, where the security check is a Bell CHSH test
rather than an error rate. Third, although the partial intercept and resend
attacker that targets only a fraction of the transmission has been studied for
BB84 and the six state protocol, it has not been applied to entanglement based
detection, and no work reports the resulting tradeoff between detection
performance and secure key rate across prepare and measure and entanglement
based protocols together.

In this work, we address these gaps through a simulation based study of 3
QKD protocols implemented in Qiskit, whose defining properties are summarized in Table~\ref{tab:protocols}. The contributions of this study are
summarized as follows.
\begin{enumerate}
\item We implement the BB84, six state, and E91 protocols within a single
Qiskit based framework, covering qubit preparation, basis reconciliation, error
reconciliation, and privacy amplification, so that the 3 protocols can be
compared on a common footing under matched ideal, synthetic noise, and real
device noise conditions.
\item Beyond the standard metrics of QBER, CHSH violation, and secure key rate,
we characterize eavesdropping detection statistically through the false positive
and false negative rates. Statistical detection for prepare and measure
protocols was established for BB84 by Lee et al. \cite{lee2022eavesdropping} and
for the six state protocol by Fiorini et al. \cite{fiorini2025sixstate}; we adopt
these thresholds and, for the first time, extend the false positive and false
negative analysis to the entanglement based CHSH detection of E91, placing
prepare and measure and entanglement based detection on a common footing.
\item We model a partial intercept and resend attacker that targets a tunable
fraction of the transmission, and we determine for each protocol the attack
strength at which detection becomes unreliable, producing a tradeoff between
detection performance and secure key rate as a function of attack strength and
noise.
\end{enumerate}

The rest of this paper is organized as follows. Section II reviews the related
work on QKD simulation, implementation, and security analysis. Section III
provides the necessary background on the BB84, six state, and E91 protocols and
the simulation environment. Section IV presents the simulation pipeline and the
performance metrics used in the evaluation. Section V presents the experimental
results and discussion.
Finally, Section VI concludes the paper and outlines directions for future
work.

\begin{table}[!ht]
\caption{Summary of the Three QKD Protocols Compared in This Work}
\label{tab:protocols}
\begin{center}
\setlength{\tabcolsep}{4pt}
\renewcommand{\arraystretch}{1.4}
\footnotesize
\begin{tabular}{|l|c|c|c|}
\hline
\textbf{Property} & \textbf{BB84} & \textbf{Six-State} & \textbf{E91} \\ \hline
Type & \multicolumn{2}{c|}{Prepare-and-measure} & Entanglement \\ \hline
Bases / settings & $Z, X$ & $Z, X, Y$ & Random axes \\ \hline
Sifting ratio & $\approx 1/2$ & $\approx 1/3$ & Matched axes \\ \hline
Full-attack signal & $25\%$ QBER & $33\%$ QBER & $S \to 1.41$ \\ \hline
Detection signal & QBER & QBER & CHSH $S$ \\ \hline
Detection threshold & $0.135$ & $0.167$ & $S = 2$ \\ \hline
\end{tabular}
\end{center}
\end{table}

\section{Related Work}
Quantum Key Distribution has attracted sustained research attention as
classical cryptography faces the prospect of quantum attacks. This section
reviews prior work on the simulation, implementation, and security analysis of
QKD protocols, with a focus on BB84 and E91. We organize the discussion into
4 themes. We first consider studies that implement both protocols together,
then studies that focus on a single protocol, then work on security analysis
and attack models, and finally simulation frameworks that support QKD research.
Table~\ref{tab:relatedwork} summarizes how these studies compare against the
present work across a set of representative features.

\subsection{Combined BB84 and E91 Implementations}
A small number of studies evaluate BB84 and E91 together, which makes them the
closest precedents for the present work. Ain et al. \cite{ain2025novel}
proposed an enhanced BB84 scheme that uses 9, 12, and 16 qubits
across 2 and 3 bases, including computational, Hadamard, and diagonal
bases, in order to strengthen security against brute force and intercept and
resend attacks. Their study reports that increasing the qubit count makes it
progressively harder for an eavesdropper to correctly guess both the
transmitted state and the measurement basis, which raises the probability of
detection. They further implemented the E91 protocol on the IBM Quantum
Experience platform and introduced a reverse gate configuration that collapses
the computational basis state from one form to another and attaches a
misleading phase angle, so that an intercepted state yields deceptive
information to the attacker. To handle the additional noise introduced by the
2 qubit gates that entanglement requires, they applied error mitigation
during the final measurement. Their results, gathered through Qiskit
simulations with 1024 shots, show fluctuations in the
measured probabilities that reveal eavesdropping and confirm that the 16
qubit configuration provides the lowest information leakage. Dhakal et al.
\cite{dhakal2025performance} provided a further direct precedent, evaluating
BB84, E91, and an enhanced BB84 variant in Qiskit under several noise models,
including depolarising, thermal relaxation, Pauli, amplitude damping, and phase
damping channels, across key lengths ranging from 200 to 3000
bits. Their study used the QBER and the CHSH violation as the primary
evaluation metrics and emphasized the trade off between security and efficiency
across the 3 protocols.

\begin{table*}[t]
\caption{Comparison of Related Work on QKD Simulation and Analysis}
\begin{center}
\setlength{\tabcolsep}{4pt}
\renewcommand{\arraystretch}{1.7}
\footnotesize
\resizebox{\textwidth}{!}{%
\begin{tabular}{|l|c|c|c|c|c|c|c|c|c|c|c|c|c|}
\hline
\textbf{Feature} &
\cite{ain2025novel} &
\cite{dhakal2025performance} &
\cite{saeed2022implementation} &
\cite{kadir2025e91} &
\cite{mariani2024simulation} &
\cite{lee2022eavesdropping} &
\cite{bechmann1999incoherent} &
\cite{singh2021role} &
\cite{diadamo2021qunetsim} &
\cite{fiorini2024estimating} &
\cite{fiorini2025sixstate} &
\cite{fiorini2026threshold} &
\textbf{Our Work} \\ \hline
Implements BB84             & \ding{51} & \ding{51} & \ding{51} & \ding{55} & \ding{55} & \ding{51} & \ding{55} & \ding{55} & \ding{55} & \ding{51} & \ding{55} & \ding{51} & \ding{51} \\ \hline
Six-state protocol          & \ding{55} & \ding{55} & \ding{55} & \ding{55} & \ding{55} & \ding{55} & \ding{51} & \ding{55} & \ding{55} & \ding{55} & \ding{51} & \ding{55} & \ding{51} \\ \hline
Implements E91              & \ding{51} & \ding{51} & \ding{55} & \ding{51} & \ding{51} & \ding{55} & \ding{55} & \ding{55} & \ding{55} & \ding{55} & \ding{55} & \ding{55} & \ding{51} \\ \hline
Noise Modeling              & \ding{51} & \ding{51} & \ding{55} & \ding{55} & \ding{51} & \ding{51} & \ding{55} & \ding{55} & \ding{51} & \ding{51} & \ding{51} & \ding{51} & \ding{51} \\ \hline
QBER / CHSH                 & \ding{51} & \ding{51} & \ding{51} & \ding{51} & \ding{51} & \ding{51} & \ding{51} & \ding{51} & \ding{55} & \ding{51} & \ding{51} & \ding{51} & \ding{51} \\ \hline
Secure Key Rate             & \ding{55} & \ding{51} & \ding{55} & \ding{51} & \ding{51} & \ding{55} & \ding{51} & \ding{51} & \ding{55} & \ding{55} & \ding{55} & \ding{55} & \ding{51} \\ \hline
FPR / FNR Detection         & \ding{55} & \ding{55} & \ding{55} & \ding{55} & \ding{55} & \ding{51} & \ding{55} & \ding{55} & \ding{55} & \ding{55} & \ding{51} & \ding{51} & \ding{51} \\ \hline
FPR / FNR for E91           & \ding{55} & \ding{55} & \ding{55} & \ding{55} & \ding{55} & \ding{55} & \ding{55} & \ding{55} & \ding{55} & \ding{55} & \ding{55} & \ding{55} & \ding{51} \\ \hline
Partial Intercept Attack    & \ding{55} & \ding{55} & \ding{55} & \ding{55} & \ding{55} & \ding{55} & \ding{55} & \ding{55} & \ding{55} & \ding{51} & \ding{51} & \ding{51} & \ding{51} \\ \hline
Multi-protocol Comparison   & \ding{51} & \ding{51} & \ding{55} & \ding{55} & \ding{55} & \ding{55} & \ding{55} & \ding{55} & \ding{55} & \ding{55} & \ding{55} & \ding{55} & \ding{51} \\ \hline
\end{tabular}%
}
\end{center}
\label{tab:relatedwork}
\end{table*}

\subsection{Single Protocol Simulations}
Several studies focus on a single protocol. Saeed et al.
\cite{saeed2022implementation} presented one of the first published quantum
circuits for BB84 in Qiskit, simulating the protocol both with and without an
eavesdropper on a local simulator and on an IBM cloud simulator. They
established the presence of an eavesdropper by observing how the probability
distributions of the measured outcomes shift when interception occurs, and they
discard and regenerate the key when this shift indicates a compromised
exchange. On the entanglement based side, Kadir and Hossain \cite{kadir2025e91}
simulated E91 on the Qiskit AerSimulator, with a third party preparing singlet
states and Alice and Bob each choosing randomly among 3 measurement bases.
They reported a CHSH parameter of 2.47, which exceeds the
classical bound of 2 and confirms a genuine Bell violation, and showed that
the key generation rate rises with the number of singlet states, reaching about
12 bits per second for 500 pairs. Mariani et al. \cite{mariani2024simulation} implemented the complete E91
protocol on the IBM Quantum platform using a depolarizing noise channel and
incorporated the classical post processing stages of error correction and
privacy amplification. In their channel model the depolarizing strength maps to
an expected QBER of 1/2 its value, and they concentrated their analysis on the
region around a 10\% error rate, near the level at which key distillation
becomes marginal. By modeling the protocol end to end, including these
post processing steps, they positioned their implementation as a reusable tool
for studying not only point to point QKD but also more complex systems that
rely on QKD for security. Badiezadegan \cite{badiezadegan2026bb84} described a
single protocol BB84 implementation that models qubits as explicit state
vectors and carries the protocol through the full classical pipeline, including
parity based block error reconciliation and binary entropy based privacy
amplification. While that work is educational and treats BB84 alone, its
post processing stages provide a concrete reference for the reconciliation and
privacy amplification steps required to compute a realistic secure key rate.

\subsection{Security Analysis and Attack Models}
A further body of work analyzes the security of QKD protocols and the attacks
to which they are exposed. Lee et al. \cite{lee2022eavesdropping} provided the
most directly relevant analysis, studying eavesdropping detection in BB84 as a
statistical decision problem. Using Hoeffding's inequality they derived upper
bounds on the false positive rate and false negative rate of detection, where
an eavesdropper is declared whenever the measured QBER reaches a threshold. By
choosing the threshold that minimizes the combined bound, an optimal value near
a QBER of $0.135$, they showed that comparing about 300 qubits keeps
both error rates below roughly 0.009\%. They also proposed a grouped BB84 variant for rapidly
varying channels. Their analysis is confined to BB84 and to an attacker that
intercepts every qubit, and they explicitly leave 3 problems open, namely
the evaluation of exact rather than upper bound error rates, the case of an
attacker that intercepts only part of the transmission, and the resulting
tradeoff between detection performance and secret key rate.

A closely related line of work by Fiorini, Garroppo, Pagano, and collaborators
develops statistical eavesdropper detection for prepare and measure protocols
under a partial intercept and resend attack, and is the closest precedent to the
prepare and measure portion of the present study. They first studied the
estimation of the interception density in BB84 using a noisy Qiskit simulator
\cite{fiorini2024estimating}, then derived two optimal QBER thresholds for the
six state protocol that minimize the false positive and false negative rates,
one from upper Hoeffding bounds and one from a limiting probability density
approach, validated on backends inspired by real IBM devices
\cite{fiorini2025sixstate}. A subsequent work proposed a BB84 threshold in which
the two communicating parties agree on the maximum fraction of key bits the
attacker may correctly intercept, removing the dependence on an estimated
interception density and setting the threshold from the system noise and the
application's security requirements \cite{fiorini2026threshold}. We adopt the
threshold construction of this line for BB84 and the six state protocol and
extend the analysis to entanglement based detection, which it does not consider.

The statistical disturbance signature that this detection exploits for the six
state protocol was established by Bechmann-Pasquinucci and Gisin
\cite{bechmann1999incoherent}, who classified the incoherent and coherent
eavesdropping strategies available against the 3 basis protocol and showed
that the additional mutually unbiased basis raises the disturbance an
intercept and resend attacker imposes relative to BB84. Their analysis provides
the theoretical basis for the larger error signal that makes six state detection
more sensitive, which the present work measures empirically as a false positive
and false negative rate. Complementing this statistical view of prepare and
measure detection, Singh et al. \cite{singh2021role} analyzed the relationship
between the Bell-CHSH violation and the secure key rate of entanglement based QKD
through a geometric representation of 2 qubit correlations. They showed that a
violation of the CHSH inequality is a necessary boundary for distilling a secret
key, though it is not on its own sufficient, since states exist that violate the
inequality yet yield no usable key, and others that yield a key only after local
filtering. Their result grounds the use of the classical bound of 2 as the
security relevant decision boundary for the CHSH based detection of E91 in the
present study, and it clarifies that crossing this boundary marks the point
beyond which the entanglement can no longer certify a secure key.
Fei et al. \cite{fei2018quantum} demonstrated a quantum man in the middle attack
that targets the calibration process of a BB84 system rather than the key
exchange itself. By intercepting calibration signals and resending faked ones,
the attacker induces a detector efficiency mismatch that opens a side channel,
and the authors proposed 2 countermeasures to remove this vulnerability. These
studies highlight that the security of QKD in practice depends not only on the
idealized protocol but also on the physical implementation and its supporting
procedures, which is relevant to any simulation that aims to model realistic
conditions.

\subsection{Simulation Frameworks}
Beyond protocol specific work, several frameworks support QKD research at the
network and circuit level. DiAdamo et al. \cite{diadamo2021qunetsim} introduced
QuNetSim, a Python framework for simulating quantum networks at the network
layer that incorporates many known quantum networking protocols and allows
users to build and test protocols, including QKD, over a range of network
configurations. Li et al. \cite{li2022pulse} extended the Quantum Toolbox in
Python with pulse level simulation of noisy circuits, enabling the modeling of
environment induced decoherence and control pulse noise that is relevant to
characterizing protocol behaviour on noisy intermediate scale quantum devices.
At the system level, Zeng et al. \cite{zeng2024towards} examined how QKD and
post quantum cryptography can be combined into hybrid schemes, weighing the
advantages and limitations of each approach and motivating the broader context
in which QKD protocols are expected to be deployed.

\subsection{Research Gap}
As summarized in Table~\ref{tab:relatedwork}, prior work has progressively built
the statistical detection picture for prepare and measure QKD but has left
entanglement based detection and the cross protocol comparison untouched. Lee et
al. \cite{lee2022eavesdropping} cast eavesdropping detection in BB84 as a
statistical decision problem and bounded the false positive and false negative
rates under a full interception attack. Fiorini et al. extended this line to a
realistic partial intercept and resend attacker, first estimating the
interception density in BB84 \cite{fiorini2024estimating}, then deriving optimal
false positive and false negative thresholds for the six state protocol under
channel noise \cite{fiorini2025sixstate}, and then proposing a refined BB84
threshold independent of the attacker's interception density
\cite{fiorini2026threshold}. These works are confined to prepare and measure
protocols and to the QBER as the detection signal. The theoretical analysis of
entanglement based security \cite{singh2021role} characterizes the CHSH
violation as a necessary boundary for a secret key but does not cast it as an
empirical detection test with measurable error rates. No prior work reports the
false positive and false negative rates of the entanglement based CHSH detector,
places prepare and measure and entanglement based detection within a single
framework under a matched partial attacker, or characterizes the resulting
tradeoff between detection performance and secure key rate across the three
protocols. The present work addresses these gaps: we adopt the established BB84
and six state thresholds, extend the false positive and false negative analysis
to the CHSH detection of E91 for the first time, and compare all three protocols
on the detection versus key rate tradeoff as the attack strength varies.

\section{System Model}
This section describes the three protocols and the threat and noise models used
throughout the study. All protocols share the same setting. Two legitimate
parties, Alice and Bob, are connected by a quantum channel that carries qubits
and by an authenticated classical channel that carries reconciliation messages.
An eavesdropper, Eve, has full access to the quantum channel and may intercept
qubits, but cannot tamper with the authenticated classical channel. This common setting is summarized in Fig.~\ref{fig:system}. Each
protocol is implemented as a quantum circuit in Qiskit and executed on the Aer
simulator, with noise introduced as described in Section~\ref{subsec:noise}.

\begin{figure}[!ht]
\centering
\includegraphics[width=\columnwidth]{Figures/fig_system_model.pdf}
\caption{System and threat model. Alice and Bob share a quantum channel carrying
qubits or entangled pairs and an authenticated classical channel for sifting,
reconciliation, and parameter estimation. Eve runs an intercept-and-resend attack on
a fraction $f$ of the quantum channel and may read but not alter the classical
channel.}
\label{fig:system}
\end{figure}

\subsection{BB84 Protocol}

In BB84, Alice generates 2 independent random bit strings, one that sets the
key bit and one that selects the preparation basis. For each position she
encodes the key bit in either the rectilinear basis $Z = \{\ket{0}, \ket{1}\}$
or the diagonal basis $X = \{\ket{+}, \ket{-}\}$. A bit value of one is encoded
by an $X$ gate, and a choice of the diagonal basis is realized by a Hadamard
gate $H$, so that the prepared state is one of the 4 signal states $\ket{0},
\ket{1}, \ket{+}, \ket{-}$, as shown in the BB84 circuit of Fig.~\ref{fig:circ_bb84}. These 4 states form 2 conjugate, or mutually
unbiased, bases, meaning that a state prepared in one basis yields a completely
random outcome when measured in the other. The protocol proceeds in 2 phases.
In the quantum phase, Alice transmits the prepared qubits one by one over the
quantum channel. Bob selects an independent random basis for each qubit,
applying a Hadamard gate before measurement when he chooses the diagonal basis
and measuring in the computational basis otherwise. In the classical phase,
carried out over the authenticated public channel, Alice and Bob announce their
basis choices, keep only the positions where the bases agree, and discard the
rest. This step, called sifting or basis reconciliation, yields a shared sifted
key. Because the two bases are chosen independently and uniformly, they agree on
about 1/2 of the positions, so the expected sifting ratio is 1/2.

The security of BB84 rests on 2 principles of quantum mechanics. The no
cloning theorem forbids an eavesdropper from copying an unknown qubit, so she
cannot keep a copy and forward the original untouched. The measurement
disturbance principle then ensures that any measurement she performs in the
wrong basis randomizes the state she resends. As a result, an eavesdropper who does
not know the preparation basis inevitably introduces discrepancies between
Alice's and Bob's sifted bits, which the parties detect by comparing a random
sample and estimating the error rate. If that rate is low enough, the remaining
sifted bits are passed to error reconciliation and privacy amplification to
produce the final secret key; if it is too high, the exchange is aborted.

\begin{figure}[!ht]
\centering
\includegraphics[width=\columnwidth]{Figures/fig_circuit_bb84.pdf}
\caption{BB84 circuit for a single qubit. Alice encodes a random bit with $X$ and selects a basis with $H$; Eve, applied with probability $f$, intercepts, measures, and resends; the channel applies depolarizing noise $\mathcal{N}_p$; Bob rotates to his random basis and measures. The $H$ gates act only for the diagonal ($X$) basis.}
\label{fig:circ_bb84}
\end{figure}

\subsection{Six-State Protocol}

The six state protocol generalizes BB84 by adding the third mutually unbiased
basis $Y = \{\ket{+i}, \ket{-i}\}$, whose states are $\ket{\pm i} = (\ket{0} \pm
i\ket{1})/\sqrt{2}$, so that Alice and Bob each select uniformly among the 3
bases $Z$, $X$, and $Y$. These 3 bases are the eigenbases of the Pauli
operators and are pairwise conjugate, which makes them the complete set of
mutually unbiased bases available to a single qubit; Fig.~\ref{fig:basis} contrasts them with the two bases of BB84. Preparation in the $Y$
basis is obtained from the computational basis by the gate sequence $S H$, and
measurement in the $Y$ basis by the inverse sequence $H S^{\dagger}$, where $S$
is the phase gate. Preparation and measurement in the $Z$ and $X$ bases are
identical to BB84. The full circuit is shown in
Fig.~\ref{fig:circ_sixstate}. All remaining steps, sifting and classical post processing,
are unchanged. Two consequences follow from the additional basis. First, because
Alice and Bob now choose independently among 3 bases, they agree on only
about 1/3 of the positions, so the sifting ratio falls from 1/2 to
1/3 and the protocol yields fewer key bits per transmitted qubit. Second,
an eavesdropper performing an intercept and resend attack must guess among 3
bases rather than 2, so she chooses correctly less often and disturbs the
state more frequently when she is wrong. The expected error she imposes on the
sifted key therefore rises from 25\% for BB84 to about 33
\% for the six state protocol. The six state protocol thus trades key
yield for a stronger and more visible response to eavesdropping, and this
explicit tension between efficiency and detectability is the property that
motivates its inclusion in this comparison.

\begin{figure}[!ht]
\centering
\includegraphics[width=0.9\columnwidth]{Figures/fig_basis_mub.pdf}
\caption{Measurement bases as axes on the Bloch sphere. BB84 uses two mutually unbiased bases, $Z$ and $X$; the six-state protocol adds the third, $Y$. The extra basis raises the intercept-and-resend disturbance from $25\%$ to $33\%$ at the cost of a smaller sifted key.}
\label{fig:basis}
\end{figure}

\begin{figure}[!ht]
\centering
\includegraphics[width=\columnwidth]{Figures/fig_circuit_sixstate.pdf}
\caption{Six-state circuit for a single qubit. The third basis $Y$ is prepared with $SH$ and measured with $S^{\dagger}H$, while the $Z$ and $X$ bases are as in BB84. Eve and the channel act as in Fig.~\ref{fig:circ_bb84}.}
\label{fig:circ_sixstate}
\end{figure}
\subsection{E91 Protocol}

The E91 protocol replaces prepared single qubits with shared entanglement and
derives its security from a Bell test rather than from basis disturbance. For
each round a source prepares a maximally entangled Bell pair. A Hadamard gate on
the first qubit, followed by a controlled NOT gate onto the second, produces the
state $(\ket{00} + \ket{11})/\sqrt{2}$. The source sends one qubit to Alice and
one to Bob, as shown in Fig.~\ref{fig:circ_e91}. Because the pair is entangled, a measurement on
one qubit instantaneously fixes the correlated outcome of the other, no matter
how far apart the 2 parties are. Alice and Bob each measure their qubit along
an independently and randomly chosen axis, implemented as a $R_y$ rotation that
selects the measurement angle followed by a computational basis measurement.
Their chosen angles fall into 2 groups. When their axes are aligned, their
outcomes are strongly correlated and the corresponding bits are kept to form the
shared key, which plays the role that the matched basis positions play in BB84.
When their axes differ in the prescribed way, the outcomes are used not for key
generation but to estimate the CHSH quantity
\begin{equation}
S = E(a, b) - E(a, b') + E(a', b) + E(a', b'),
\label{eq:chsh}
\end{equation}
where $a, a'$ and $b, b'$ are Alice's and Bob's alternative measurement axes and
$E(\cdot, \cdot)$ is the correlation between their outcomes for a given pair of
axes, defined as the average of the product of the $\pm 1$ outcomes. For a
genuinely entangled pair measured at the optimal angles, $S$ can reach the
Tsirelson bound $2\sqrt{2} \approx 2.83$, whereas any classical or local hidden
variable description obeys the bound $S \leq 2$. A measured value of $S$ above
2 therefore certifies that genuine entanglement survived transmission. An
eavesdropper who intercepts and measures a qubit destroys part of this
entanglement, which weakens the correlations and drives $S$ back toward the
classical bound, so the collapse of $S$ toward two is the signature that the
channel has been disturbed. This entanglement based check is appealing because
it tests the physical resource that secures the key directly, rather than
inferring an attack from an error rate alone.

\begin{figure}[!ht]
\centering
\includegraphics[width=\columnwidth]{Figures/fig_circuit_e91.pdf}
\caption{E91 circuit. A Bell pair is created with $H$ and CNOT and shared between Alice and Bob; the channel applies $\mathcal{N}_p$; Eve, with probability $f$, measures Bob's qubit along a random axis and resends; Alice and Bob measure along randomly chosen axes $R_y(a)$ and $R_y(b)$. Matched axes form the key and the remaining settings estimate $S$.}
\label{fig:circ_e91}
\end{figure}

\subsection{Threat Model: Partial Intercept and Resend}
\label{subsec:threat}
We adopt the intercept and resend attack as the adversary model. For an attacked
qubit, Eve measures in a randomly chosen basis and forwards a freshly prepared
qubit consistent with her outcome, which disturbs the state whenever her basis
differs from the preparation basis. Following the partial intercept and resend
model introduced for BB84 by Fiorini et al. \cite{fiorini2024estimating} and
applied by them to the six state protocol \cite{fiorini2025sixstate}, we
parameterize the attack by an interception fraction $f \in [0, 1]$, so that Eve
attacks each qubit, or each entangled pair, independently with probability $f$.
The case $f = 1$ recovers the full attack studied in prior work, while
intermediate values model an adversary that limits her interference to remain
below the detection threshold. We apply this attacker for the first time to the
entanglement based E91 protocol, so that a single adversary model can be compared
across all three protocols. For E91, interception of one member of a pair
collapses the entanglement and reduces the measured CHSH value.



\subsection{Noise Model}
\label{subsec:noise}
Two noise regimes are considered. The first is a synthetic per qubit depolarizing
channel of strength $p$, applied uniformly to every transmitted qubit so that the
modeled disturbance does not depend on which gates a particular qubit happened to
receive. This regime provides a clean, controllable axis along which to vary the
channel quality. The second is a device noise model extracted from the
\texttt{ibm\_marrakesh} backend. Unlike the uniform depolarizing channel, this
model reproduces the heterogeneous error sources of real hardware, including per
gate error rates, thermal relaxation characterized by the qubit $T_1$ and $T_2$
times, and readout errors that vary from qubit to qubit. Testing against these
physical realities matters because they do not behave as a single tunable
parameter; instead they natively inflate the baseline QBER and degrade the
baseline CHSH value $S$ even when no eavesdropper is present. Because honest
sessions therefore begin closer to the abort condition, the device model
inherently shrinks the stealth zone, the band of attack strengths an eavesdropper
can occupy before the protocol is forced to abort, and so provides a stricter and
more realistic test of detection than synthetic noise alone. The 2 regimes are
treated as mutually exclusive, so that any reported error or detection statistic
can be attributed to a single, clearly specified noise source.

\subsection{Classical Post-Processing}
For BB84 and the six state protocol, the sifted key is processed in 3 stages.
Parameter estimation reveals a random subset of the sifted bits to measure the
QBER, error reconciliation corrects the surviving mismatches using parity based
block correction, and privacy amplification compresses the reconciled key to
remove the information potentially leaked to Eve, with the compression length
set from the binary entropy of the estimated QBER and the number of disclosed
parity bits. For E91, the correlated matched basis outcomes furnish the key,
which is then processed through the same reconciliation and privacy
amplification stages, while the unmatched settings supply the CHSH estimate.
\begin{figure}[t]
\centering
\includegraphics[width=0.95\columnwidth]{Figures/fig_pipeline.pdf}
\caption{End-to-end simulation pipeline. Each run proceeds from circuit construction, through one of the two mutually exclusive noise regimes (synthetic depolarizing or the \texttt{ibm\_marrakesh} device model), the partial intercept-and-resend eavesdropper acting on a fraction $f$ of the transmission, and the classical post-processing stages, to the final computation of the error indicators and detection statistics.}
\label{fig:pipeline}
\end{figure}
\section{Performance Metrics}
\label{sec:metrics}
The end to end simulation pipeline is summarized in Fig.~\ref{fig:pipeline}. Each
run begins with circuit construction for the chosen protocol, after which noise is
injected from one of the two regimes of Section~\ref{subsec:noise}, either the
synthetic depolarizing channel or the \texttt{ibm\_marrakesh} device model. The
partial intercept and resend eavesdropper of Section~\ref{subsec:threat} then acts
on a fraction $f$ of the transmission. The surviving signals pass through sifting,
error reconciliation, and privacy amplification, and the run concludes by
computing the error indicators and detection statistics, namely the QBER or CHSH
value, the secure key rate, and the false positive and false negative rates.



We evaluate the three protocols using 2 groups of metrics. The first group
contains standard quantities that validate the implementation and place it
alongside prior work. The second group contains the detection statistics that
form the contribution of this study.

\subsection{Basic Metrics}
\textit{Quantum Bit Error Rate (QBER).} For BB84 and the six state protocol, the
QBER is the fraction of positions in the sifted key where Alice's and Bob's bits
disagree. It captures the combined effect of channel noise and eavesdropping and
is the primary error indicator for prepare and measure protocols.

\textit{CHSH parameter.} For E91, the CHSH value $S$ of \eqref{eq:chsh} plays the
analogous role. A value above the classical bound of 2 certifies entanglement,
and the amount by which $S$ exceeds 2 measures the quality of the correlations
that secure the key.

\textit{Secure key rate.} For all three protocols, the secure key rate is the
number of secret key bits produced per transmitted qubit. Following the
asymptotic Shor and Preskill expression, the secret fraction is taken as $r =
\max(0,\, 1 - 2 H_2(Q))$, where $Q$ is the QBER and $H_2$ is the binary entropy
function, and the secure key rate is the product of this fraction and the
sifting ratio. The rate falls to 0 once the QBER reaches the threshold near
11\%, beyond which no secret key can be distilled.

\textit{Sifting efficiency.} The sifting ratio is the fraction of transmitted
qubits retained after basis reconciliation. It is about 1/2 for BB84 and
about 1/3 for the six state protocol, and it quantifies the key yield cost
of adding a third basis.

\subsection{Novel Metrics: Statistical Detection}
\textit{False positive rate (FPR).} Eavesdropping is declared when the measured
error indicator crosses a decision threshold, the QBER for BB84 and the six
state protocol, and a corresponding threshold on $S$ for E91. The false positive
rate is the probability that this test declares an eavesdropper when none is
present and only channel noise is acting. A high FPR means a protocol aborts
honest sessions unnecessarily.

\textit{False negative rate (FNR).} The false negative rate is the probability
that the test fails to declare an eavesdropper when one is in fact present. A
high FNR means an attacker passes undetected. The FPR and FNR together
characterize the reliability of eavesdropping detection, and we estimate both
empirically from repeated simulation runs rather than from analytic bounds. Fig.~\ref{fig:detection} shows how both error rates arise from the overlap of the honest and eavesdropper distributions at the decision threshold. For BB84 and
the six state protocol the statistical detection framework follows Lee et al.
\cite{lee2022eavesdropping} and Fiorini et al. \cite{fiorini2025sixstate}; the
contribution here is to report the same false positive and false negative rates
for the entanglement based CHSH detection of E91, which has not previously been
analyzed in these terms, and to compare all three on a common footing.

\begin{figure*}[ht]
\centering
\includegraphics[width=\textwidth]{Figures/fig_detection_decision.pdf}
\caption{Statistical definition of the detection errors. An alarm is declared when
the error indicator crosses a fixed threshold. (a) For BB84 and the six-state
protocol the alarm fires when the QBER reaches the threshold, so the false positive
rate is the honest-channel mass above it and the false negative rate is the
eavesdropper mass below it. (b) For E91 the alarm fires when the CHSH value $S$ falls
below the classical bound $S=2$, mirroring the roles. Both rates are estimated
empirically by Monte Carlo.}
\label{fig:detection}
\end{figure*}

\textit{Decision thresholds.} Two different limits involve the QBER, and we keep
them strictly separate. The first is the secure key rate limit near a QBER of
11\%, the asymptotic Shor and Preskill bound beyond which the distilled
secret fraction reaches 0 and no key can be safely produced. This limit only
dictates whether a key can still be distilled; it does not raise an alarm. The
second, used for detection, is the statistical threshold that actually triggers
the eavesdropper alarm, and it is a separate quantity. We initially explored
calibrating the detection thresholds dynamically from each channel's baseline
honest noise profile, placing the alarm a fixed margin above the observed
honest level in order to suppress false alarms on noisy channels. We ultimately
discarded this adaptive scheme in favor of fixed, literature grounded limits, so
that the reported error rates are measured against rigorous and standardized
security bounds rather than boundaries fitted to our own data. For BB84 we adopt
the optimal QBER threshold of Lee et al. \cite{lee2022eavesdropping}, who place
the decision boundary that minimizes the combined Hoeffding bound on the 2 error
rates at a QBER of $0.135$. For the six state protocol we adopt the Hoeffding
balanced threshold derived by Fiorini et al. \cite{fiorini2025sixstate}. When the
QBER is estimated from $n$ compared bits, Hoeffding's inequality bounds the two
detection errors by
\begin{equation}
\mathrm{FPR} \le e^{-2n(\tau - q_0)^2}, \qquad
\mathrm{FNR} \le e^{-2n(D - \tau)^2},
\label{eq:hoeffding}
\end{equation}
where $\tau$ is the decision threshold, $q_0$ is the honest error rate, and $D$ is
the expected disturbance under a full intercept and resend attack. Equalizing the
two exponents places the threshold at the midpoint $\tau = (q_0 + D)/2$. For BB84,
with $D = 1/4$, this coincides with the optimized value of Lee et al.
\cite{lee2022eavesdropping}. For the six state protocol the full attack
disturbance is larger, $D = 1/3$, giving $\tau = 1/6 \approx 0.167$ with a
noiseless reference $q_0 = 0$, consistent with Fiorini et al. We use these
established thresholds for the prepare and measure protocols and contribute the
corresponding analysis for the entanglement based detector of E91, for which no
statistical detection threshold has previously been reported. For E91 the
decision boundary is the CHSH classical bound of $2.0$, the value used to
separate secure from compromised channels in prior entanglement based studies
\cite{ekert1991quantum, dhakal2025performance, kadir2025e91}. A Bell violation is
a necessary condition for distilling a secret key \cite{singh2021role}, so a
measured $S$ that falls to 2 signals that the entanglement, and with it the
security of the channel, can no longer be certified.

\textit{Detection resource cost.} For a target FPR and FNR, we report the number
of compared qubits, or entangled pairs, that each protocol requires to reach
that target, which quantifies the quantum resources detection consumes.

\textit{Detection versus key rate tradeoff.} Finally, by sweeping the
interception fraction $f$ of Section~\ref{subsec:threat}, we report for each
protocol the joint behaviour of detection performance and secure key rate, and
we identify the attack strength at which detection becomes unreliable. This
tradeoff is the central result the framework is designed to produce.

\section{Results and Discussion}
\label{sec:results}
Unless otherwise stated, every synthetic data point is the mean over $50$
independent Monte Carlo trials, using $1000$ transmitted qubits per trial for
BB84 and the six state protocol and $8192$ entangled pairs per trial for E91; the
device noise points use $30$ trials, with $200$ qubits or $1024$ pairs per trial.
The interception fraction $f$ is swept from $0$ to $1$ in steps of $0.1$ for the
synthetic channels and in steps of $0.25$ for the device model. We first validate
the implementation against known analytic limits, then present the detection
signal as a function of attack strength, the stealth window and the resulting
tradeoff against secure key rate, the robustness of each detector to channel
noise, the behaviour under real device noise, and finally the resource cost of
detection.

\begin{figure*}[t]
\centering
\includegraphics[width=0.95\textwidth]{Figures/fig_metric_vs_f.pdf}
\caption{Detection signal versus attack strength on an ideal channel. (a) The
QBER of BB84 and the six-state protocol rises monotonically with the interception
fraction $f$ and reaches the fixed detection thresholds ($0.135$ and $0.167$) only
near $f = 0.5$. (b) The E91 CHSH value $S$ falls from the Tsirelson bound toward
the classical bound $S = 2$, which it crosses near $f = 0.6$. Shaded bands show
one standard deviation over the Monte Carlo trials.}
\label{fig:metric_f}
\end{figure*}

\subsection{Implementation Validation}
On an ideal channel with no eavesdropper, the 3 protocols reproduce their
textbook limits, which confirms that the circuits and the post processing are
correct. BB84 yields a QBER of $0.000$ and a secure key rate of $0.498$, matching
the expected 1/2 sifting ratio, and the six state protocol yields a QBER of
$0.000$ and a secure key rate of $0.333$, matching its 1/3 sifting ratio.
For E91 the honest CHSH value is $S = 2.823$, within statistical error of the
Tsirelson bound $2\sqrt{2} \approx 2.83$, and its secure key rate is $0.223$,
consistent with the 2/9 setting combinations used for key generation.
Under a full intercept and resend attack ($f = 1$) the disturbance signatures
also match theory: the BB84 QBER rises to $0.253$ and the six state QBER to
$0.338$, close to the predicted 25 and 33\%, while the
E91 CHSH value falls to $S = 1.415 \approx \sqrt{2}$, well inside the classical
region. In every case the secure key rate is 0 under the full attack. These
anchor values appear as the endpoints of Fig.~\ref{fig:metric_f}.

\subsection{Detection Signal Versus Attack Strength}
Fig.~\ref{fig:metric_f} shows how each protocol's error indicator responds as the
interception fraction grows on an otherwise ideal channel. The QBER of BB84 and
the six state protocol rises monotonically with $f$, and the CHSH value of E91
falls monotonically, so each indicator moves steadily from its honest value
toward its full attack value. The rate of change differs across protocols. The
six state QBER rises fastest, reaching its detection threshold of $0.167$ between
$f = 0.4$ and $f = 0.5$, because its 3 basis measurement makes Eve guess
wrong more often. The BB84 QBER reaches its lower threshold of $0.135$ between
$f = 0.5$ and $f = 0.6$, and the E91 CHSH value crosses the classical bound
$S = 2$ between $f = 0.5$ and $f = 0.6$. The important consequence is that for all
3 protocols the indicator does not reach its decision boundary until roughly
1/2 of the transmission is intercepted, which already suggests that a weaker
attacker can remain below the alarm.

\begin{figure}[t]
\centering
\includegraphics[width=\columnwidth]{Figures/fig_tradeoff.pdf}
\caption{The tradeoff between detection and key rate on an ideal channel. (a) The
false negative rate stays at one through $f = 0.3$ (shaded stealth window) and
then falls, earliest for the six-state protocol and latest for E91. (b) The secure
key rate declines with $f$ and is already near zero by the attack strength at
which each protocol's false negative rate first drops below one half.}
\label{fig:tradeoff}
\end{figure}

\subsection{Stealth Window and Detection Versus Key Rate Tradeoff}
Fig.~\ref{fig:tradeoff} makes this tension explicit by plotting, against the same
attack fraction, the false negative rate of detection in the upper panel and the
secure key rate in the lower panel. Three regimes are visible. For $f \le 0.3$,
shaded in the figure, every protocol has a false negative rate of 1, meaning a
single session never raises the alarm, yet the secure key rate is still clearly
positive. At $f = 0.1$ the secure key rates are $0.322$, $0.191$, and $0.158$ for
BB84, the six state protocol, and E91 respectively, and they remain positive
through $f = 0.3$. This band is a genuine stealth window, in which the attacker
harvests usable key while staying invisible to per session detection. As $f$
increases further the false negative rate falls, but the order in which the
protocols become detectable matters. The six state protocol is the first to
respond, with its false negative rate dropping to $0.42$ at $f = 0.5$, followed by
BB84, which falls to $0.20$ at $f = 0.6$, and finally E91, which is still at
$0.54$ at $f = 0.6$ and only reaches 0 at $f = 0.7$. The six state protocol
therefore closes the stealth window earliest and E91 sustains it longest, which
makes the entanglement based detector the easiest for a partial attacker to
evade. The lower panel adds the decisive point: by the time each protocol detects
the attacker reliably, its secure key rate has already collapsed to near zero. The
secure key rate is below $0.01$ for every protocol at the attack strength where
its false negative rate first drops below 1/2. In other words, within a
single session at these sample sizes, the fixed thresholds catch the attacker
essentially only once the key is already destroyed, so detecting a stealthy low
fraction attacker is not possible from one exchange and instead requires comparing
many more signals, which we quantify in Section~\ref{subsec:resource}.

\begin{figure}[t]
\centering
\includegraphics[width=\columnwidth]{Figures/fig_fpr_vs_noise.pdf}
\caption{False positive rate on an honest channel as the synthetic noise level
increases. BB84 and the six-state protocol remain near zero, whereas the E91 CHSH
test raises a false alarm on most honest sessions once the channel error rate
exceeds about 8\%.}
\label{fig:fpr_noise}
\end{figure}

\subsection{Robustness to Channel Noise}
A detector is only useful if it does not raise the alarm on an honest but noisy
channel. Fig.~\ref{fig:fpr_noise} reports the false positive rate of each protocol
as the synthetic channel noise increases with no eavesdropper present. BB84 and
the six state protocol are robust: their false positive rate stays at 0 up to
an honest error rate of 8\%, and only BB84 shows a small rate of $0.08$
at the 11\% level, because its lower threshold of $0.135$ sits closer to
the honest error rate than the six state threshold of $0.167$. E91 behaves very
differently. Its honest CHSH value degrades steadily with noise, from $2.823$ at
0 noise to $2.611$, $2.291$, $1.989$, and $1.726$ at the 2, 5, 8, and
11\% levels, so it drops below the classical bound near the 8
\% level. As a result the E91 false positive rate jumps to $0.60$ at 8
\% and to $1.00$ at 11\%, meaning the CHSH test declares an
eavesdropper on essentially every honest session once the channel is moderately
noisy. This fragility is the principal weakness of entanglement based detection in
this study, and it reflects the fact that the CHSH value is a more delicate
quantity than the QBER.

\begin{figure*}[t]
\centering
\includegraphics[width=0.95\textwidth]{Figures/fig_device.pdf}
\caption{Behaviour under the \texttt{ibm\_marrakesh} device noise model. (a) The
honest secure key rate falls only moderately relative to the ideal channel, and no
protocol raises a false alarm. (b) The false negative rate versus interception
fraction follows the same shape as on the ideal channel, so the stealth window
persists on real hardware.}
\label{fig:device}
\end{figure*}

\subsection{Behaviour Under Device Noise}
The synthetic results above isolate a single noise strength, whereas a real device
mixes many error sources. Fig.~\ref{fig:device} reports the 3 protocols under
the noise model extracted from the \texttt{ibm\_marrakesh} backend. The honest
results are reassuring. As the left panel shows, real device noise reduces the
secure key rate only moderately, from $0.498$ to $0.414$ for BB84, from $0.333$ to
$0.287$ for the six state protocol, and from $0.223$ to $0.166$ for E91, and the
false positive rate is 0 for all 3. In particular the honest E91 CHSH value
on the device is $S = 2.697$, comfortably above the classical bound, because the
effective error rate of this backend is below the 8\% level at which the
synthetic CHSH test broke down. This is an important qualification of the
fragility result: E91 is fragile to large channel noise, but the measured device
noise is mild enough that honest entanglement based operation remains viable. The
right panel shows that detection on the device follows the same pattern as on the
ideal channel, with the false negative rate staying near 1 for a 1/4
interception and falling through 1/2 near $f = 0.5$ for all 3 protocols,
so the stealth window persists on real hardware.

\begin{figure*}[t]
\centering
\includegraphics[width=0.95\textwidth]{Figures/fig_resource_cost.pdf}
\caption{Detection resource cost. (a) False positive rate versus the number of
compared signals $K$ on a worst-case honest noisy channel (11\% error
for BB84 and the six-state protocol, 5\% for E91). (b) False negative
rate versus $K$ against a full attacker on a clean channel. The six-state protocol
suppresses false alarms with the fewest samples, and the false negative rate is
easy to drive down for all three protocols.}
\label{fig:resource}
\end{figure*}

\subsection{Detection Resource Cost}
\label{subsec:resource}
Because a single session cannot catch a stealthy attacker, the relevant question
is how many compared signals are needed to make detection statistically reliable.
Fig.~\ref{fig:resource} reports the false positive rate and the false negative
rate as a function of the number of compared signals $K$, where the false positive
rate is measured on a worst case honest noisy channel and the false negative rate
on a fully intercepted clean channel. To keep each protocol near its own marginal
honest regime, the honest channel used for the false positive measurement is held
at an 11\% error rate for BB84 and the six state protocol and at a 5
\% error rate for E91, since E91 at the 11\% level alarms on every
honest run regardless of $K$. The false negative rate is easy to drive down for
all 3 protocols, reaching 5\% or below within a few hundred compared
signals. The binding constraint is the false positive rate. The six state protocol
is the most sample efficient, falling below 5\% by about $K = 300$, BB84
requires about $K = 1000$, and E91 lies in between, falling below 5\%
between $K = 500$ and $K = 1000$. All 3 reach a negligible false positive rate
within a few thousand signals. The ranking mirrors the threshold geometry: the six
state threshold of $0.167$ leaves the widest margin above an 11\% honest
channel, so it suppresses false alarms with the fewest samples, whereas the BB84
threshold of $0.135$ leaves the least margin and is the most sample hungry.

\subsection{Summary of Findings}

\begin{table*}[t]
\caption{Summary of the Comparative Detection Results Across the Three Protocols}
\label{tab:summary}
\begin{center}
\setlength{\tabcolsep}{6pt}
\renewcommand{\arraystretch}{1.3}
\footnotesize
\begin{tabular}{|l|c|c|c|c|c|c|}
\hline
\textbf{Protocol} & \textbf{Detection} & \textbf{Stealth-window} & \textbf{Detection onset} & \textbf{FPR at} & \textbf{$K$ for} & \textbf{Device honest} \\
 & \textbf{threshold} & \textbf{edge ($f$)} & \textbf{($f$, FNR $\le 0.5$)} & \textbf{$11\%$ noise} & \textbf{FPR $\le 5\%$} & \textbf{SKR} \\ \hline
BB84      & QBER $0.135$ & $0.3$ & $0.6$ & $0.08$ & $\approx 1000$ & $0.414$ \\ \hline
Six-State & QBER $0.167$ & $0.3$ & $0.5$ & $0.00$ & $\approx 300$  & $0.287$ \\ \hline
E91       & $S = 2.0$    & $0.5$ & $0.7$ & $1.00$ & $\approx 700$  & $0.166$ \\ \hline
\end{tabular}
\end{center}
\end{table*}

Table~\ref{tab:summary} consolidates the comparison across the three protocols.
Taken together, the results identify the six state protocol as offering the most
favourable detection profile among the 3. It produces the strongest
disturbance signal per intercepted qubit, closes the stealth window at the lowest
attack fraction, never raises a false alarm across the synthetic noise range
studied, and needs the fewest compared signals to make detection reliable, at the
cost of the lowest sifted key yield. BB84 is intermediate on every axis. E91
directly certifies the entangled resource that secures its key, but it is the most
evadable by a partial attacker and by far the most fragile to channel noise,
although real device noise is mild enough that honest operation is preserved.
Across all 3 protocols the central finding is the same: a partial intercept
and resend attacker has a stealth window at low interception fractions in which
key still flows and per session detection fails, and reliable detection of such an
attacker is a question of accumulated statistics rather than of any single
exchange.

\section{Conclusion}
We presented a unified, noise aware Qiskit study of the BB84, six state, and E91
protocols that evaluates eavesdropping detection statistically, through the false
positive and false negative rates, and extends this analysis from full
interception to a partial intercept and resend attacker that targets a tunable
fraction of the transmission. Statistical detection of this kind was established
for BB84 by Lee et al. \cite{lee2022eavesdropping} and recently extended to the
six state protocol by Fiorini et al. \cite{fiorini2025sixstate}; building on those
thresholds, we provide the first false positive and false negative analysis of
the entanglement based CHSH detector and the first comparison of statistical
detection across prepare and measure and entanglement based protocols within one
framework. Across matched ideal, synthetic noise, and real device conditions, the
study produces for each protocol a tradeoff between detection performance and
secure key rate as a function of attack strength. The main findings are that a
partial attacker has a stealth window at low interception fractions in which a
usable key is still produced while a single session never raises the alarm, that
reliable per session detection appears only once the attack is strong enough to
have already destroyed the key, and that detecting a stealthy attacker is
therefore a matter of accumulated statistics. Among the 3 protocols the six state
protocol offers the most favourable detection profile, closing the stealth window
earliest, never raising a false alarm across the noise range studied, and
requiring the fewest compared signals, at the cost of the lowest key yield; BB84
is intermediate; and E91 directly certifies entanglement but is the most evadable
and the most fragile to channel noise, although honest operation is preserved
under the measured device noise.

Several directions remain open. The E91 detection boundary used here is the fixed
CHSH classical bound $S = 2$, a security limit rather than a statistically
optimized threshold; deriving the threshold on $S$ that balances the false
positive and false negative rates from the variance of the CHSH estimator, in the
way Lee et al. \cite{lee2022eavesdropping} and Fiorini et al.
\cite{fiorini2025sixstate} do for the QBER, would give the entanglement based
detector the same statistical footing as the prepare and measure detectors. The
six state detection threshold was placed by balancing the Hoeffding bounds on the
two detection errors; a tighter construction that instead minimizes their weighted
sum, together with distribution aware rather than worst case bounds, would refine
it further. The secure key rate used the asymptotic Shor and Preskill expression,
and a finite key treatment would tighten the resource cost results. Because the
central finding is that detection is a matter of accumulated statistics, a
sequential detector that aggregates evidence across sessions is a natural way to
catch a low fraction attacker that no single session catches. Finally, executing
the protocols on real quantum hardware rather than a device noise model, and
extending the adversary beyond intercept and resend to collective and coherent
attacks, would test whether the stealth window and the protocol ordering observed
here persist under stronger conditions.

\section*{Acknowledgment}
The authors thank Mitacs for the Globalink Research Award, which provided the
opportunity to carry out this research, and the University of Aberdeen for
hosting the research visit during which this work was conducted.

\bibliographystyle{ieeetr}
\bibliography{references}

\end{document}